\documentclass[11pt]{article}
\usepackage{fontspec,amsmath,amssymb,geometry}
\setmainfont{Times New Roman}\geometry{margin=0.7in}
\newcommand{\topic}[1]{\par\leftskip=0pt\section*{\normalfont\large #1}}\newcommand{\initial}[1]{{\Large #1}}\newcommand{\result}[1]{\par\leftskip=0pt\medskip\noindent\hspace*{1em}{\bfseries #1}\par\nobreak\leftskip=1em\noindent}\renewcommand{\textbf}[1]{\result{#1}}
\newenvironment{chaptermap}{\begin{center}\begin{minipage}{0.82\textwidth}\small}{\end{minipage}\end{center}\medskip}
\title{Solving Linear Equations}\author{T.T.}\date{}
\begin{document}\maketitle
\begin{chaptermap}\itshape A system is interpreted by rows and columns. Elimination produces echelon form, pivots determine existence and uniqueness, and matrix factorizations organize the same computation.\end{chaptermap}
\topic{\initial{S}YSTEMS AND \initial{E}LIMINATION}
\[Ax=b,\qquad Ax=x_1a_1+\cdots+x_na_n.\]
\textbf{Row operations:} interchange rows; multiply a row by a nonzero constant; add a multiple of one row to another.
\textbf{Rank criterion:} $Ax=b$ is solvable iff $\operatorname{rank}(A)=\operatorname{rank}([A\ b])$.
\textbf{Pivot theorem:} free variables occur in non-pivot columns; if $r$ is the rank, the solution set has $n-r$ free parameters.
\[Ax=0\Rightarrow x=x_{\rm pivot}+x_{\rm free};\quad x=x_p+x_n,\quad Ax_n=0.\]
\textbf{Rouché--Capelli conclusion:} a consistent system has either one solution (no free variables) or infinitely many solutions.
\topic{\initial{M}ATRIX \initial{O}PERATIONS}
\[A^{-1}A=I,\quad (AB)^{-1}=B^{-1}A^{-1},\quad (A^T)^{-1}=(A^{-1})^T.\]
\[E_{ij}^{-1}=E_{ij}(-c),\quad E_i(c)^{-1}=E_i(c^{-1}),\quad P^{-1}=P^T.\]
\textbf{Elimination factorization:} $EA=U$, hence $A=LU$; with row exchanges, $PA=LU$.
\textbf{LU solution:} solve $Ly=b$ followed by $Ux=y$.
\topic{\initial{S}PECIAL \initial{M}ATRICES}
\[A^TA\text{ is symmetric},\quad A=S+K,\quad S=\frac{A+A^T}{2},\quad K=\frac{A-A^T}{2}.\]
\textbf{Orthogonal matrix:} $Q^TQ=I\Longleftrightarrow Q^{-1}=Q^T$; it preserves dot products and lengths.
\[\det(AB)=\det A\det B,\qquad A^{-1}=\frac{1}{\det A}\operatorname{Cof}(A)^T.\]
\end{document}
